\section{Conclusion and Recommendations}

Three TO-BE scenarios were evaluated against the AS-IS baseline using paired t-tests at the 5\% significance level. The combined scenario yields the greatest improvement, reducing mean cycle time by 12.11 minutes relative to AS-IS, while also producing the narrowest confidence interval, indicating a more consistent process flow. Scheduling alone reduces mean cycle time by 6.71 minutes and pre-processing by 5.47 minutes, confirming that both interventions contribute independently and that their effects are additive when combined. Based on these results, Dr. Weems is advised to pursue both interventions simultaneously. The scheduling optimization only requires administrative coordination, making this a natural first step. Pre-processing reduces load on Attending but requires organizational commitment from residents and attending physicians. As total patients processed remains the same in these scenarios, the improvements reflect reduced waiting time rather than increased throughput. This suggests that further gains would require more direct intervention at the Attending bottleneck. 